\documentclass[10pt,a4paper]{amsart}
\usepackage[margin=19mm]{geometry}
\usepackage{amsmath,amssymb,mathtools}
\usepackage[scr=rsfs,cal=euler]{mathalfa}
\usepackage{xfrac}
\usepackage[unicode,hidelinks,bookmarks=false]{hyperref}
\newcommand{\ie}{i.e.\ }
\newcommand{\cf}{cf.\ }
\newcommand{\resp}{respectively\ }
\newcommand{\Subsection}{\subsection}
\providecommand{\nobreakdash}{\nobreak\hbox{-}}
\setlength{\emergencystretch}{2em}
\multlinegap0pt
\usepackage{bm}
\usepackage{bbm}
\usepackage{dsfont}
\usepackage{xspace}
\usepackage{cancel}
\usepackage{mathtools}
\usepackage{amsthm}
\makeatletter
\@ifundefined{theorem}{\newtheorem{theorem}{Theorem}}{}
\@ifundefined{lemma}{\newtheorem{lemma}[theorem]{Lemma}}{}
\@ifundefined{proposition}{\newtheorem{proposition}{Proposition}}{}
\makeatother
\title[Irregularity and Topological Indices in Fibonacci Word Trees]{Irregularity and Topological Indices in Fibonacci Word Trees and Modified Fibonacci Word Index}
\author{Jasem Hamoud, Duaa Abdullah}
\date{Source: arXiv:2505.00602v1 (CC BY 4.0)}
\begin{document}
\maketitle

\noindent\emph{Source and licence.} Irregularity and Topological Indices in Fibonacci Word Trees and Modified Fibonacci Word Index. Version arXiv:2505.00602v1 (CC BY 4.0), distributed under the Creative Commons Attribution 4.0 International licence, as recorded in the arXiv licence metadata for this exact version and shown on the abstract page https://arxiv.org/abs/2505.00602v1. Full rights notice: licenses/arxiv-2505.00602-CC-BY-4.0.txt.\par\smallskip

\noindent\textit{Editorial scope.} Counted unit: the Lemma whose source label is \texttt{lemn1} in the article (source0.tex lines 399--405, source label \texttt{lemn1}), delivered together with the article's own complete original proof of it (source0.tex lines 406--449, one \texttt{proof} environment of 44 lines). No mathematics is written, repaired or re-derived: statement and proof are the article's own text line for line. Required context retained uncounted: pron1 (lines 189--191); pron2 (lines 385--387). Every definition, display and label of those lines is retained unchanged. Omitted from this package and not redistributed: 1--188, 192--384, 388--398, 450--912. Nothing inside a retained excerpt is omitted or altered: every retained body is delivered exactly as the article's lines are written, so no omission receipt of any kind accompanies this package. Citations: the retained excerpts carry no citation command at all, so no bibliography entry of the article is transcribed and no reference is redistributed. Reference adaptation: none was needed and none was made; every reference inside the delivered unit resolves inside it, so no unresolved reference or citation is produced. Printed slips kept literal: no printed word, symbol, hypothesis or number of the counted statement or of its proof was altered. Layout and class: the article's own class is \texttt{article} with packages including graphicx, amssymb, stmaryrd, enumitem, geometry, tikz, float, url, mathtools, amsmath, amsfonts, amsthm, mathrsfs, mma; the wrapper is the lane's pinned \texttt{amsart} editorial wrapper at 10pt on A4, which restates the theorem-like environments the retained text needs and adds the article's own macro definitions that the retained text needs (none), with extra wrapper packages (bm, bbm, dsfont, xspace, cancel, mathtools). The delivered numbering is the unit's own. Assets: no retained span contains a figure environment, a graphics-inclusion command, a PDF-page inclusion, a table or a TikZ picture; no image file is copied into this package, the package declares no asset dependency and no image file is redistributed with it. Licence: version arXiv:2505.00602v1 of this article is distributed under the Creative Commons Attribution 4.0 International licence, as recorded in the arXiv licence metadata for this exact version and shown on the abstract page https://arxiv.org/abs/2505.00602v1. A full rights notice is delivered at licenses/arxiv-2505.00602-CC-BY-4.0.txt. Characters: every retained excerpt is pure ASCII, so the delivered engine, XeLaTeX with its default Latin Modern text font, reports no missing character.
\begin{proposition}~\label{pron1}
Let $\mathscr{T}$ be a tree with maximum degree $\Delta$ then $\operatorname{FWI}^*(\mathscr{T})\ge \Delta(\Delta^2 - 1)$ with equality if and only if $\mathscr{T}$ is isomorphic to either a path or a tree containing only one vertex of $\deg (V(\mathscr{T}))>2$.
\end{proposition}

\begin{proposition}~\label{pron2}
Let $P_n(\mathscr{T})$ be the path of a tree $\mathscr{T}$ of $n$-vertex where $n\geq 5$, let $S_n(\mathscr{T})$ be the star of $\mathscr{T}$, then $\operatorname{FWI}^* (P_n) < \operatorname{FWI}^*(\mathscr{T}) <\operatorname{FWI}^*(S_n)$.
\end{proposition}

\begin{lemma}~\label{lemn1}
Let $\mathscr{T}$ be a tree of Fibonacci word, let $\alpha$ and $\beta$ are non-adjacent vertices in $\mathscr{T}$ where $\deg(F_\alpha\ge \deg (F_\beta)$ then
\[
\operatorname{FWI}^*(\mathscr{T})=\operatorname{FWI}^*(\mathscr{T}-\alpha\beta)+2[(2\deg(F_\alpha)+1)r_\alpha+(2\deg(F_\beta)+1)r_\beta]-\Gamma(\mathscr{T}),
\]
where $\Gamma(\mathscr{T})=3\deg(F_\alpha)(\deg (F_\alpha)+1)-  \deg(F_\beta)(\deg (F_\beta)+1)$.
\end{lemma}

\begin{proof}
Firstly, we provided a motivation to stabilise the Fibonacci word by Proposition~\ref{pron1},\ref{pron2} as $\operatorname{FWI}^*(\mathscr{T})\ge \Delta(\Delta^2 - 1)$, then we have:  

\begin{align*}
\operatorname{FWI}^*(\mathscr{T}-\alpha\beta) - \operatorname{FWI}^*(\mathscr{T})
&=(\deg(F_\alpha)+1)^2 - (\deg(F_\beta)+1)^2 \\
& + \sum_{x\in N(\alpha)} \bigg( \big|(\deg(F_\alpha)+1)^2 - (\deg(F_x))^2\big| - \big|(\deg(F_\alpha))^2 - (\deg(F_x))^2\big|\bigg)\\
& + \sum_{y\in N(\beta)} \bigg( \big|(\deg(F_\beta)+1)^2 - (\deg(F_y))^2\big| - \big|(\deg(F_\beta))^2 - (\deg(F_y))^2\big|\bigg)\,.
\end{align*}
Now, consider $N(\alpha)=L(\alpha)\cup E(\alpha)\cup R(\alpha)$. We aim to determine the union of the sets $ L(\alpha) $, $ E(\alpha) $, and $ R(\alpha) $, defined as follows:
\[
\begin{cases}
L(\alpha) = \{\beta \in N(\alpha) : \deg(F_\beta) < \deg(F_\alpha)\}, \\
E(\alpha) = \{\beta \in N(\alpha) : \deg(F_\beta) = \deg(F_\alpha)\}, \\
R(\alpha) = \{\beta \in N(\alpha) : \deg(F_\beta) > \deg(F_\alpha)\},
\end{cases}
\]
where $ N(\alpha) $ is a set containing elements $ \beta $, and $ F_\beta $ and $ F_\alpha $ are polynomials associated with elements $ \beta $ and $ \alpha $, respectively, with $ \deg(\cdot) $ denoting the degree of a polynomial. The objective is to compute $ L(\alpha) \cup E(\alpha) \cup R(\alpha) $. The sets $ L(\alpha) $, $ E(\alpha) $, and $ R(\alpha) $ are subsets of $ N(\alpha) $, defined based on the degree of the polynomial $ F_\beta $ relative to the degree of $ F_\alpha $. Specifically:
\begin{itemize}
    \item $ L(\alpha) $ contains all $ \beta \in N(\alpha) $ such that $ \deg(F_\beta) < \deg(F_\alpha) $.
    \item $ E(\alpha) $ contains all $ \beta \in N(\alpha) $ such that $ \deg(F_\beta) = \deg(F_\alpha) $.
    \item $ R(\alpha) $ contains all $ \beta \in N(\alpha) $ such that $ \deg(F_\beta) > \deg(F_\alpha) $.
\end{itemize}
For any $ \beta \in N(\alpha) $, the value of $ \deg(F_\beta) $ must satisfy exactly one of the following mutually exclusive conditions relative to $ \deg(F_\alpha) $:
\begin{enumerate}
    \item $ \deg(F_\beta) < \deg(F_\alpha) $, placing $ \beta $ in $ L(\alpha) $,
    \item $ \deg(F_\beta) = \deg(F_\alpha) $, placing $ \beta $ in $ E(\alpha) $,
    \item $ \deg(F_\beta) > \deg(F_\alpha) $, placing $ \beta $ in $ R(\alpha) $.
\end{enumerate}
Same results if we consider $N(\alpha)=L(\beta)\cup E(\beta)\cup R(\beta)$, then we have: 

\begin{align*}
\operatorname{FWI}^*(\mathscr{T}-\alpha\beta) - \operatorname{FWI}^*(\mathscr{T})
&=(\deg F_\alpha - \deg F_\beta)(\deg F_\alpha +\deg F_\beta +2)+ (2\deg F_\alpha+1)(e_\alpha +l_\alpha - r_\alpha)\\
& \quad + (2\deg F_\beta+1)(e_\beta +l_\beta - r_\beta)\,,
\end{align*}
which is equivalent to
\begin{align*}
\operatorname{FWI}^*(\mathscr{T}-\alpha\beta) - \operatorname{FWI}^*(\mathscr{T})
&=3\deg F_\alpha(\deg F_\alpha+1) + \deg F_\beta(\deg F_\beta-1)\\
&-2[(2\deg F_\alpha+1)r_\alpha + (2\deg F_\beta+1)r_\beta]\,.
\end{align*}
As desire.
\end{proof}

\end{document}
